\section{Smaller positive systems and circuit evaluations}
\label{sec:low-degree}

The coefficient adapter also supports two improvements that preserve the
positive-domain arithmetic constructions of package~30 \cite{repo30}.  One
retains degree four and reduces the arithmetic evaluation cost; the other
retains twelve unknowns and reduces the cost of its much higher-degree norm.
All the identities in this section hold with the four coefficient ports
treated as independent indeterminates.  We then specialize \emph{all four}
ports to the new complete $121$ by $8801$ tuple.  The resulting fixed
polynomials need not equal the historical fixed polynomials, because their
coefficients have changed.  Their arithmetic templates have the exact
identities proved below.

\begin{center}
\begin{tabular}{lrrrr}
\hline
Presentation & Unknowns & Equations & Degree & Operations\\
\hline
Quadratic comparison system &31&22&2&60\\
Its sum of squares &31&1&4&125\\
Twelve-variable product &12&1&1081&108\\
Twelve-variable system &12&3&1072&103\\
Transport-shifted system &12&3&1072&102\\
\hline
\end{tabular}
\end{center}

Every unknown in this table is a strictly positive integer.  For a system,
``degree'' means the maximum degree of a residual.  Every binary addition,
subtraction and multiplication is charged, including a multiplication by a
fixed numeral; a comparison and reuse of an existing register are free.
The source costs at the same two points were $65/130$ operations for the
quadratic system/quartic, and $107/111$ for the twelve-variable
system/product.  These are local improvements.  The separate inherited
$59$-operation higher-degree systems and $67$-operation single equation
retain their lower total-operation counts.

\subsection{A 60-operation quadratic system}

The $31$ supplied positive integers are
\[
\begin{gathered}
J,F,\alpha,z,f,h,i,o,j,s,w,\tau,\eta,\zeta,y,\gamma,\\
N,E,P,c,a,X,D,C_2,L,F_2,V,g,I,B_{\rm aux},\Delta.
\end{gathered}
\]
Use the fixed ports $\beta,K_0,M_C,M_F$, where $M_F$ is the already shifted
high mask.  Introduce only the polynomial abbreviations
\[
\begin{gathered}
q=\beta J+1,\quad k=\eta+\zeta,\quad H=4a+3,
\quad C=q-F-\alpha,\quad U=L+1,\\
Y=a-E,\qquad Q=c-\eta,\qquad T=P-k,\qquad R=jc-V.
\end{gathered}
\]
In particular, $R$ is quadratic.  The following are the $22$ comparisons:
\begin{equation}
\begin{aligned}
 U&=q^2,&N&=qU,&X&=wq,&Y&=sN,\\
 E&=XY,&Q&=kY,&T&=EQ,&\tau^2&=PT+1,\\
 D&=X+ac+\gamma H,&\Delta&=a^2+H,&C_2&=c^2,\\
 D^2&=\Delta C_2+1,&k&=R+1+hE.
\end{aligned}
\label{eq:quadratic-first}
\end{equation}
\begin{equation}
\begin{aligned}
 R&=(q-F)(N-q)-CL+(M_C+qM_F)J,\\
 (K_0+w)C&=F+z\beta J,\\
 F_2&=f^2,\qquad V=of-c,\qquad g=y^2-V^2,\\
 I&=iC_2,\qquad B_{\rm aux}=I^2,
       \qquad B_{\rm aux}=\Delta(F_2-1),\\
 y^2-B_{\rm aux}g&=1.
\end{aligned}
\label{eq:quadratic-second}
\end{equation}
Each equality has a quadratic residual after substituting the displayed
abbreviations.  If these residuals are $r_1,\ldots,r_{22}$ in the displayed
order, the single polynomial is
\begin{equation}
 \boxed{\qquad Q_{125}=\sum_{j=1}^{22}r_j^2.\qquad}
 \label{eq:q125}
\end{equation}
Thus its positive zero set is exactly the system's positive solution set.

The first four simplifications reorient existing equalities so that a
computed difference or sum serves two purposes.  The equality
$L=q^2-1$ becomes $L+1=q^2$, and the same $L+1$ occurs in $N=q(L+1)$.
The equality $a=E+sN$ becomes $a-E=sN$, reusing the $a-E$ in $E=X(a-E)$.
Similarly, $c-\eta=k(a-E)$ reuses $c-\eta$ in the equation defining $P$,
and $P-k=E(c-\eta)$ reuses $P-k$ in $\tau^2=P(P-k)+1$.
These are polynomial identities on all integer assignments.

The fifth simplification is the positive transport-coordinate shift also
used in package~26 \cite{repo26}.  Its original equation is
\[
 (K_0+w)C+(q-F)-z_{\rm old}(q-1)=1.
\]
Substituting $z_{\rm old}=z+1$ and $q-1=\beta J$ gives exactly the second
equation in \eqref{eq:quadratic-second}.  A new positive solution therefore
gives an old one.  Conversely, first apply the established compiler
reconstruction to an old solution.  It gives
\[
 C>0,\qquad F<q,\qquad X=wq=2^R,\qquad R>3q,
\]
and hence $w>q>F$.  The original transport equation implies
\[
 (z_{\rm old}-1)(q-1)=(K_0+w)C-F>0.
\]
Thus $z_{\rm old}\ge2$, and $z=z_{\rm old}-1$ remains positive.  This
establishes a bijection of positive zero sets.  The positivity claim is a
consequence of the completed source reconstruction; it is not presumed in
the initial bootstrap.

The resulting comparison circuit has $33$ multiplications and $27$
additions/subtractions, hence $60$ operations.  Forming the residuals,
squaring them and summing gives $55$ multiplications and $70$
additions/subtractions, hence $125$ operations.  Its degree is exactly
four: $\tau$ occurs only in the residual
$\tau^2-P(P-k)-1$, so the coefficient of $\tau^4$ in \eqref{eq:q125}
is exactly one.  This argument is uniform in every fixed coefficient
slice and does not depend on a generic numerical specialization.

The supplied builder checks exact equality of all $22$ residuals with
the source residuals after $z_{\rm old}=z+1$, treating all four ports
symbolically.  It expands the quartic to $299$ terms when port exponents
are included in the symbolic monomials.  Specializing actual coefficients
can identify terms, so $299$ is stated with that convention.  The
degree-four certificate is unaffected.

\subsection{A 108-operation twelve-variable equation}

The twelve supplied positive integers are now
\[
 J,F,\alpha,z,h,i,T,s,w,\eta,\zeta,t.
\]
The notation in this subsection is local.  Set
\begin{align*}
 q&=\beta J+1,& C_w&=q-F-\alpha,&X&=wq,&Y&=sq^3,&E&=XY,\\
 k&=\eta+\zeta,&c&=kY+\eta,&a&=E+Y,&H&=4a+3,&\Delta&=a^2+H,\\
 \ell&=EkY,&A&=1+\ell(\ell+k),&G&=1+\Delta c^2,&b_0&=X+ac.
\end{align*}
The transport and index registers are
\begin{align*}
 R&=((q-1)(q-F)+\alpha)(q^2-1)+(M_C+qM_F)J,\\
 N_k&=k-hE-R,\qquad
 N_t=(K_0+w)C_w+q-F-z(q-1).
\end{align*}
For the Pell norm define
\[
\begin{gathered}
 I=ic^2,\quad D=1+\Delta I^2,\quad Q=\Delta(D-1),\quad d=Q-1,
 \quad b=c+RD,\\
 M=k(I^2+T^2),\qquad L=HM+1,\qquad W=b_0+Ht.
\end{gathered}
\]
The economical scalar evaluator is
\begin{align*}
 S&=H^2Q,\qquad H_P=SDc^2T^2,\qquad H_U=Sb^2,\\
 n&=(HM)^2A,\qquad d_n=dn,\qquad H_e=H_P+H_U-H^2,\\
 j&=4H_PH_U,\qquad \mathcal T=2d(W^2+L^2G),\\
 r&=d(W^2-L^2G),\qquad\mathcal K=r^2,\\
 a_1&=\mathcal T-2(d_n+H_e),\qquad b_1=(d_n-H_e)^2-j-\mathcal K,\\
 v&=a_1(a_1\mathcal T+2b_1)-16d_nj,\\
 u&=a_1^2\mathcal K-b_1^2.
\end{align*}
All these letters denote polynomial abbreviations, not new witnesses.
Finally put
\begin{equation}
 N=\mathcal K v^2-u(\mathcal T v-u),\qquad
 \boxed{\quad P_{108}=N_kN_t(1+N)-1.\quad}
 \label{eq:p108}
\end{equation}
Its equation is $P_{108}=0$.  Its three-equation version is
\begin{equation}
 N_k=1,\qquad N_t=1,\qquad
 \mathcal K v^2=u(\mathcal T v-u).
 \label{eq:system103}
\end{equation}
Replacing the middle equation by
\begin{equation}
 (K_0+w)C_w=F+z(q-1)
 \label{eq:positive102transport}
\end{equation}
gives the $102$-operation system.  Its $z$ is the old transport coordinate
minus one; the positivity and equivalence proof is exactly the one above,
applied after the twelve-variable reconstruction.

\subsection{Why the compressed norm is exact}

The source scalar norm is the product over the main square-root conjugates
of a polynomial encoding the first, strong and auxiliary Pell conditions.
For its quadratic-remainder evaluator, the two final coefficients are
\[
 v_r=a_1(a_1\mathcal T+2b_1)-16d_nj,
 \qquad u_r=b_1^2-a_1^2\mathcal K,
\]
and the final norm is
\[
 u_r^2+\mathcal T u_rv_r+\mathcal K v_r^2.
\]
The identity
\begin{equation}
 u_r^2+\mathcal T u_rv_r+\mathcal K v_r^2
 =\mathcal K v^2-u(\mathcal T v-u),\qquad
 u=-u_r,\quad v=v_r,
 \label{eq:factored-positive-norm}
\end{equation}
holds in every commutative ring.  The sign change reverses an existing
subtraction, and the factorization saves one multiplication.

Two additional ring identities give the other savings.  Sharing $H^2Q$
between $H_P=H^2QDc^2T^2$ and $H_U=H^2Qb^2$ saves one multiplication.
Computing
\[
 d^2(W^2-L^2G)^2=[d(W^2-L^2G)]^2
\]
saves another because the source's separate $d^2$ has no other consumer.
Thus the product takes $108=65+43$ operations, the summands being
multiplications and additions/subtractions.  The final comparison in
\eqref{eq:system103} avoids one more subtraction, giving
$103=63+40$.  The transport shift gives $102=63+39$.

No new rational-root, root-separation or Pell-saturation hypothesis is
introduced by these changes.  The norm is the identical four-port
polynomial template.  Consequently the source's universal reconstruction,
its product-unit sign argument and its nonzero leading homogeneous form
remain valid.  The norm has exact degree $1072$, and $N_k,N_t$ have degrees
$7,2$, respectively, so \eqref{eq:p108} has exact degree $1081$.
The active coefficient tuple has $\beta>0$, which is the coefficient
condition needed for the inherited noncancellation certificate.

\subsection{Logical consequence and reproducibility}

For each presentation in this section, the new complete coefficient
adapter, the inherited arithmetic reconstruction, and the literal K215
calibration give
\begin{align*}
 \text{positive solvability}
 &\quad\Longleftrightarrow\quad \text{blank halting of K215}\\
 &\quad\Longleftrightarrow\quad \text{existence of a finite ZFC refutation}.
\end{align*}
Negation yields the requested equivalence with $\operatorname{Con}(\mathrm{ZFC})$
inside ZFC.  The construction establishes neither side of that equivalence.

The quartic generator performs exact sparse-polynomial verification of
all residuals and of the degree certificate.  The twelve-variable checker
verifies five local ring identities exactly, audits every live input and
operation, compares $240$ signed modular assignments with the source
template, and checks twelve additional assignments using residues of the
actual bound coefficients.  Both active bindings authenticate the same
four payloads.  Their historical source files have been restored to exact
repository bytes and checked against the pinned Git blob identifiers.
The new algebraic and coefficient arguments remain ordinary mathematical
proofs; these finite checks are not a new end-to-end Lean or Rocq theorem.
